\chapter{留数}

\section{孤立奇点与留数的定义}
我们称复平面上的点 \( z_0 \) 是函数 \( f(z) \) 的**孤立奇点**，当且仅当 \( f(z) \) 在 \( z_0 \) 的一个去心邻域 \( B^\circ(z_0, \delta) \) 上解析。

设 \( z_0 \) 是 \( f(z) \) 的孤立奇点，其洛朗展开为
\[
f(z) = \sum_{n=-\infty}^{+\infty} a_n (z - z_0)^n,
\]
则 \( \frac{1}{z - z_0} \) 项的系数 \( a_{-1} \) 称为 \( f(z) \) 在 \( z_0 \) 处的**留数**，记为 \( \text{Res}[f(z), z_0] \)。

对 \( \delta < R \)，有
\[
\oint_{|z - z_0| = \delta} f(z) dz = 2\pi i \, \text{Res}[f(z), z_0].
\]

\subsection{留数基本定理}
设 \( C \) 是简单光滑闭曲线，\( f(z) \) 在 \( C \) 所围闭区域上除内部点 \( a_1, a_2, \dots, a_n \) 外解析，则
\[
\oint_C f(z) dz = 2\pi i \sum_{k=1}^n \text{Res}[f(z), a_k].
\]

对有界光滑连通闭区域 \( D \)，有
\[
\int_{\partial D} f(z) dz = 2\pi i \sum_{k=1}^n \text{Res}[f(z), a_k].
\]

\section{孤立奇点的分类}
设 \( z_0 \) 是 \( f(z) \) 的孤立奇点，其洛朗展开为 \( \sum_{n=-\infty}^{+\infty} a_n (z - z_0)^n \)：
\begin{itemize}[leftmargin=2em]
    \item 负部分为零：\( z_0 \) 是**可去奇点**；
    \item 负部分为有限项：\( z_0 \) 是**极点**；
    \item 负部分为无限项：\( z_0 \) 是**本性奇点**。
\end{itemize}

\( z_0 \) 是可去奇点 \( \iff \lim_{z \to z_0} f(z) \) 存在有限；
\( z_0 \) 是极点 \( \iff \lim_{z \to z_0} f(z) = \infty \)；
\( z_0 \) 是本性奇点 \( \iff \lim_{z \to z_0} f(z) \)（广义意义下）不存在。

\subsection{零点与极点的关系}
\( z_0 \) 是 \( f(z) \) 的 \( m \) 阶零点 \( \iff z_0 \) 是 \( \frac{1}{f(z)} \) 的 \( m \) 阶极点。

设 \( f(z) = \frac{P(z)}{Q(z)} \)，\( z_0 \) 是 \( P(z) \) 的 \( m \) 阶零点、\( Q(z) \) 的 \( n \) 阶零点：
\begin{itemize}[leftmargin=2em]
    \item \( m > n \)：\( z_0 \) 是 \( f(z) \) 的 \( m-n \) 阶零点；
    \item \( m = n \)：\( z_0 \) 是 \( f(z) \) 的可去奇点；
    \item \( m < n \)：\( z_0 \) 是 \( f(z) \) 的 \( n-m \) 阶极点。
\end{itemize}

\subsection{Picard 大定理}
设 \( z_0 \) 是 \( f(z) \) 的本性奇点，\( f(z) \) 在去心邻域 \( B^\circ(z_0, \delta) \) 上解析，则 \( \mathbb{C} \setminus f(B^\circ(z_0, \delta)) \) 中至多有一个复数。

\section{留数的计算}
\begin{enumerate}[leftmargin=2em]
    \item \textbf{线性性}：
    \[
    \begin{split}
    \text{Res}\left[f(z) \pm g(z), z_0\right] &= \text{Res}\left[f(z), z_0\right] \pm \text{Res}\left[g(z), z_0\right], \\
    \text{Res}\left[a f(z), z_0\right] &= a \text{Res}\left[f(z), z_0\right], \quad a \in \mathbb{C}.
    \end{split}
    \]
    \item 若 \( z_0 \) 是可去奇点，则 \( \text{Res}\left[f(z), z_0\right] = 0 \)。
    \item 若 \( z_0 \) 是 \( m \) 阶极点，则
    \[
    \text{Res}\left[f(z), z_0\right] = \frac{1}{(m-1)!} \lim_{z \to z_0} \frac{d^{m-1}}{dz^{m-1}} \left[(z - z_0)^m f(z)\right].
    \]
    \item 若 \( f(z) = \frac{g(z)}{h(z)} \)，\( z_0 \) 是 \( h(z) \) 的一阶零点且不是 \( g(z) \) 的零点，则
    \[
    \text{Res}\left[f(z), z_0\right] = \frac{g(z_0)}{h'(z_0)}.
    \]
\end{enumerate}

\section{无穷远点的留数}
无穷远点 \( \infty \) 是 \( f(z) \) 的孤立奇点，当且仅当存在 \( R>0 \)，使得 \( f(z) \) 在 \( |z|>R \) 上解析。

\( \infty \) 是 \( f(z) \) 的可去奇点（\( m \) 阶极点、本性奇点）\( \iff 0 \) 是 \( g(\xi) = f\left( \frac{1}{\xi} \right) \) 的可去奇点（\( m \) 阶极点、本性奇点）。

留数定义：
\[
\text{Res}[f(z), \infty] := -\text{Res}\left[ f\left( \frac{1}{\xi} \right) \frac{1}{\xi^2}, 0 \right].
\]

若 \( f(z) \) 在扩充复平面上只有有限多个奇点 \( \{a_1, \dots, a_n, \infty\} \)，则
\[
\sum_{k=1}^{n} \text{Res}\left[f(z), a_k\right] + \text{Res}\left[f(z), \infty\right] = 0.
\]